\documentclass[11pt]{article}
\usepackage{amsmath,amssymb,amsthm}
\usepackage{geometry}
\geometry{margin=1in}
\usepackage{hyperref}
\hypersetup{colorlinks,linkcolor=blue,citecolor=blue,urlcolor=blue}
\newtheorem{theorem}{Theorem}[section]
\newtheorem{lemma}{Lemma}[section]
\newtheorem{assumption}{Assumption}[section]
\newtheorem{remark}{Remark}[section]

\begin{document}

\title{Frank–Wolfe (Conditional Gradient) for Smooth Non‑Convex Optimization \\ Detailed Algorithm, Theory, and Complete Proof}
\author{}
\date{}
\maketitle

\tableofcontents
\newpage

%%%%%%%%%%%%%%%%%%%%%%%%%%%%%%%%%%%%%%%%%%%%%%%%%%%%%%%%%%%%
\section*{Algorithm Name}
\addcontentsline{toc}{section}{Algorithm Name}
\begin{center}
\Large \textbf{Frank–Wolfe (Conditional Gradient) Method}
\end{center}
%%%%%%%%%%%%%%%%%%%%%%%%%%%%%%%%%%%%%%%%%%%%%%%%%%%%%%%%%%%%

\section{Mathematical Formulation}
\label{sec:formulation}
Let 
\[
f:\mathbb{R}^{d}\rightarrow\mathbb{R}
\]
be a differentiable (possibly non‑convex) function and let 
\[
\mathcal C\subset\mathbb{R}^{d}
\]
be a non‑empty compact convex set (the feasible domain).  
Denote by 
\[
\nabla f(x) \in\mathbb{R}^{d}
\]
the gradient of $f$ at $x$.  

At iteration $k\ge 0$ the algorithm performs the following two steps:

\begin{enumerate}
    \item \textbf{Linear minimization oracle (LMO)}  
    \[
    s_{k}\;:=\;\arg\min_{s\in\mathcal C}\;\langle \nabla f(x_{k}),\,s\rangle .
    \tag{1}
    \]
    \item \textbf{Convex‑combination update} with a stepsize $\gamma_{k}\in[0,1]$  
    \[
    x_{k+1}\;:=\;(1-\gamma_{k})\,x_{k}\;+\;\gamma_{k}\,s_{k}.
    \tag{2}
    \]
\end{enumerate}

The quantity
\[
\mathcal G(x_{k})\;:=\;\max_{s\in\mathcal C}\langle \nabla f(x_{k}),\,x_{k}-s\rangle
    \;=\;-\min_{s\in\mathcal C}\langle \nabla f(x_{k}),\,s\rangle
\tag{3}
\]
is called the \emph{Frank–Wolfe (FW) gap}. Because $s_{k}$ solves (1), we have
\[
\langle \nabla f(x_{k}),\,x_{k}-s_{k}\rangle \;=\;\mathcal G(x_{k}).
\tag{4}
\]

%%%%%%%%%%%%%%%%%%%%%%%%%%%%%%%%%%%%%%%%%%%%%%%%%%%%%%%%%%%%
\section{Pseudocode}
\label{sec:pseudocode}
\begin{verbatim}
Input :   initial point x0 ∈ C
          stepsize rule {γk}k≥0 ⊂ [0,1]
          maximum iterations T
Output:   sequence {xk}k=0..T

x ← x0
for k = 0,…,T‑1 do
    // 1) Linear minimization oracle
    sk ← argmin_{s∈C} <∇f(x), s>

    // 2) Stepsize (choose one of the rules below)
    //    (a) Fixed:          γk = γ ∈ (0,1]
    //    (b) Diminishing:    γk = 2/(k+2)
    //    (c) Adaptive:       γk = min { G(x)/(L D^2), 1 }
    //       where G(x) = <∇f(x), x‑sk> and L is the smoothness constant

    // 3) Convex‑combination update
    x ← (1‑γk) * x + γk * sk
end for
return {xk}
\end{verbatim}
All quantities used in the stepsize rules are defined in Section~\ref{sec:theory}.

%%%%%%%%%%%%%%%%%%%%%%%%%%%%%%%%%%%%%%%%%%%%%%%%%%%%%%%%%%%%
\section{Dry Run Example}
\label{sec:dryrun}
Consider the one‑dimensional problem
\[
f(w)=(w-3)^{2},
\qquad \mathcal C=[0,\,5],
\qquad w_{0}=0,
\qquad \gamma_{k}=0.1\;\;(\text{fixed}).
\]

The gradient is $\nabla f(w)=2(w-3)$.  The LMO solves
\[
s_{k}= \arg\min_{s\in[0,5]} \; 2(w_{k}-3)\,s .
\]
Because the objective is linear in $s$, the minimizer is an endpoint:
\[
s_{k}= 
\begin{cases}
5, & \text{if } w_{k}<3,\\[2mm]
0, & \text{if } w_{k}>3.
\end{cases}
\]

\begin{center}
\begin{tabular}{c|c|c|c|c}
$k$ & $w_{k}$ & $\nabla f(w_{k})$ & $s_{k}$ & $w_{k+1}=0.9\,w_{k}+0.1\,s_{k}$ \\ \hline
0 & $0$ & $-6$ & $5$ & $0.5$ \\
1 & $0.5$ & $-5$ & $5$ & $0.95$ \\
2 & $0.95$ & $-4.1$ & $5$ & $1.355$ 
\end{tabular}
\end{center}

The corresponding objective values are
\[
f(w_{0})=9,\quad f(w_{1})=(0.5-3)^{2}=6.25,\quad
f(w_{2})=(0.95-3)^{2}=4.2025,\quad
f(w_{3})=(1.355-3)^{2}=2.709\,.
\]

The FW gap at iteration $k$ equals $\mathcal G(w_{k})=|\nabla f(w_{k})|\;|w_{k}-s_{k}|$,
e.g. $\mathcal G(w_{0})=6\cdot5=30$, $\mathcal G(w_{1})=5\cdot4.5=22.5$, etc.,
showing a monotone reduction.

%%%%%%%%%%%%%%%%%%%%%%%%%%%%%%%%%%%%%%%%%%%%%%%%%%%%%%%%%%%%
\section{Assumptions}
\label{sec:assumptions}
\begin{assumption}[Compact feasible set]
\label{ass:C}
The set $\mathcal C\subset\mathbb{R}^{d}$ is convex, compact, and has finite
diameter
\[
D \;:=\; \max_{x,y\in\mathcal C}\|x-y\| < \infty .
\]
\end{assumption}

\begin{assumption}[Smoothness of the objective]
\label{ass:L}
The function $f$ is differentiable and its gradient is $L$‑Lipschitz continuous,
i.e.
\[
\|\nabla f(x)-\nabla f(y)\|\le L\|x-y\|,\qquad\forall x,y\in\mathbb{R}^{d}.
\tag{5}
\]
Equivalently,
\[
f(y)\le f(x)+\langle\nabla f(x),y-x\rangle+\frac{L}{2}\|y-x\|^{2},
\qquad\forall x,y.
\tag{6}
\]
\end{assumption}

\begin{assumption}[Existence of a finite optimal value]
\label{ass:opt}
Define $f^{*}:=\inf_{x\in\mathcal C}f(x)$.  Because $\mathcal C$ is compact and $f$
is continuous, $f^{*}>-\infty$.
\end{assumption}

No convexity of $f$ is required; the analysis holds for arbitrary smooth
(non‑convex) objectives.

%%%%%%%%%%%%%%%%%%%%%%%%%%%%%%%%%%%%%%%%%%%%%%%%%%%%%%%%%%%%
\section{Main Convergence Theorem}
\label{sec:theorem}
\begin{theorem}[Convergence of Frank–Wolfe for smooth non‑convex $f$]
\label{thm:main}
Assume \ref{ass:C}–\ref{ass:opt}.  
Let $\{x_{k}\}_{k\ge0}$ be generated by the algorithm with the
\emph{adaptive stepsize}
\[
\gamma_{k}\;:=\;\min\!\Bigl\{\frac{\mathcal G(x_{k})}{L D^{2}},\,1\Bigr\},
\tag{7}
\]
where $\mathcal G(x_{k})$ is the FW gap defined in \eqref{3}.  
Then for any integer $T\ge 1$,
\[
\boxed{\;
\min_{0\le k\le T-1}\mathcal G(x_{k})
\;\le\;
\sqrt{\frac{2L D^{2}\bigl(f(x_{0})-f^{*}\bigr)}{T}}
\;}
\tag{8}
\]
Consequently, the algorithm produces an $\varepsilon$‑stationary point
(i.e. a point with $\mathcal G(x)\le\varepsilon$) after at most
\[
T = \Bigl\lceil \frac{2 L D^{2}\bigl(f(x_{0})-f^{*}\bigr)}{\varepsilon^{2}}\Bigr\rceil
\]
iterations.
\end{theorem}

\begin{remark}
When $f$ is convex, $\mathcal G(x)$ upper‑bounds the primal gap
$f(x)-f^{*}$, and inequality \eqref{8} recovers the classic $O(1/T)$ rate for
the Frank–Wolfe method with line‑search or the $O(1/\sqrt{T})$ rate for the
fixed stepsize $\gamma_{k}=2/(k+2)$.  In the non‑convex case the bound is
expressed in terms of the FW gap, which serves as a first‑order stationarity
measure for constrained problems.
\end{remark}

%%%%%%%%%%%%%%%%%%%%%%%%%%%%%%%%%%%%%%%%%%%%%%%%%%%%%%%%%%%%
\section{Theoretical Properties}
\label{sec:properties}
\begin{enumerate}
\item \textbf{Stability.}  
   The update \eqref{2} is a convex combination of two points inside
   $\mathcal C$, hence $x_{k}\in\mathcal C$ for all $k$ (invariance of the
   feasible set).

\item \textbf{Stepsize sensitivity.}  
   The adaptive rule \eqref{7} automatically scales the step according to the
   local curvature $L$ and the size of the feasible region $D$.  If
   $\mathcal G(x_{k})$ is small the algorithm takes a tiny step, preventing
   overshooting; if the gap is large it may take a full step ($\gamma_{k}=1$).

\item \textbf{Comparison to baseline methods.}  
   Compared with projected gradient descent (PGD), Frank–Wolfe avoids an
   explicit Euclidean projection; the per‑iteration cost is the LMO,
   which can be substantially cheaper for structured $\mathcal C$ (e.g.
   simplex, trace‑norm ball).  The convergence rate \eqref{8} matches the
   $O(1/\sqrt{T})$ guarantee of PGD for non‑convex smooth problems when
   measured by the same stationarity criterion.
\end{enumerate}

%%%%%%%%%%%%%%%%%%%%%%%%%%%%%%%%%%%%%%%%%%%%%%%%%%%%%%%%%%%%
\section{Preliminaries (Definitions and Lemmas)}
\label{sec:prelim}
\begin{lemma}[Descent Lemma from $L$‑smoothness]
\label{lem:descent}
For any $x,y\in\mathbb{R}^{d}$,
\[
f(y) \le f(x) + \langle\nabla f(x),y-x\rangle + \frac{L}{2}\|y-x\|^{2}.
\tag{9}
\]
\end{lemma}
\begin{proof}
Inequality \eqref{9} is exactly the smoothness inequality \eqref{6}.
\end{proof}

\begin{lemma}[Bound on the step length]
\label{lem:stepbound}
For any $k$, the distance moved by the algorithm satisfies
\[
\|x_{k+1}-x_{k}\| = \gamma_{k}\,\|s_{k}-x_{k}\|\le \gamma_{k} D .
\tag{10}
\]
\end{lemma}
\begin{proof}
Since $s_{k},x_{k}\in\mathcal C$ and $\operatorname{diam}(\mathcal C)=D$,
$\|s_{k}-x_{k}\|\le D$.  Multiplying by $\gamma_{k}$ yields \eqref{10}.
\end{proof}

%%%%%%%%%%%%%%%%%%%%%%%%%%%%%%%%%%%%%%%%%%%%%%%%%%%%%%%%%%%%
\section{Main Proof (Step‑by‑Step Derivation)}
\label{sec:proof}
We now prove Theorem~\ref{thm:main}.  All non‑trivial manipulations are
annotated with a step number of the form \textbf{[Step N]}.

\subsection*{Step 1 – Apply the descent lemma}
Using Lemma~\ref{lem:descent} with $x=x_{k}$ and $y=x_{k+1}$,
\[
f(x_{k+1}) \le f(x_{k}) + \langle\nabla f(x_{k}),x_{k+1}-x_{k}\rangle
               + \frac{L}{2}\|x_{k+1}-x_{k}\|^{2}.
\tag{11}
\]

\subsection*{Step 2 – Substitute the update rule}
From the algorithm \eqref{2},
\[
x_{k+1}-x_{k}= \gamma_{k}(s_{k}-x_{k}).
\tag{12}
\]
Insert \eqref{12} into \eqref{11}:
\[
f(x_{k+1}) \le f(x_{k}) + \gamma_{k}\langle\nabla f(x_{k}),s_{k}-x_{k}\rangle
               + \frac{L}{2}\gamma_{k}^{2}\|s_{k}-x_{k}\|^{2}.
\tag{13}
\]

\subsection*{Step 3 – Relate the inner product to the FW gap}
By definition of the FW gap \eqref{3} and the optimality of $s_{k}$ in the
LMO \eqref{1},
\[
\langle\nabla f(x_{k}),s_{k}-x_{k}\rangle = -\mathcal G(x_{k}).
\tag{14}
\]
Thus \eqref{13} becomes
\[
f(x_{k+1}) \le f(x_{k}) - \gamma_{k}\,\mathcal G(x_{k})
               + \frac{L}{2}\gamma_{k}^{2}\|s_{k}-x_{k}\|^{2}.
\tag{15}
\]

\subsection*{Step 4 – Upper‑bound the squared distance}
By Lemma~\ref{lem:stepbound}, $\|s_{k}-x_{k}\|\le D$, therefore
\[
\frac{L}{2}\gamma_{k}^{2}\|s_{k}-x_{k}\|^{2}
   \le \frac{L D^{2}}{2}\,\gamma_{k}^{2}.
\tag{16}
\]

\subsection*{Step 5 – Combine the inequalities}
From \eqref{15}–\eqref{16} we obtain
\[
f(x_{k+1}) \le f(x_{k}) - \gamma_{k}\,\mathcal G(x_{k})
               + \frac{L D^{2}}{2}\,\gamma_{k}^{2}.
\tag{17}
\]

\subsection*{Step 6 – Choose the adaptive stepsize}
Recall the stepsize definition \eqref{7}:
\[
\gamma_{k}= \min\!\Bigl\{\frac{\mathcal G(x_{k})}{L D^{2}},\,1\Bigr\}.
\tag{18}
\]
We consider the two possible cases separately.

\paragraph{Case A: $\displaystyle \gamma_{k}= \frac{\mathcal G(x_{k})}{L D^{2}}\le 1$}
Plugging $\gamma_{k}$ into \eqref{17} yields
\[
\begin{aligned}
f(x_{k+1})
&\le f(x_{k})
   - \frac{\mathcal G(x_{k})^{2}}{L D^{2}}
   + \frac{L D^{2}}{2}\,\Bigl(\frac{\mathcal G(x_{k})}{L D^{2}}\Bigr)^{2} \\[1mm]
&= f(x_{k}) - \frac{\mathcal G(x_{k})^{2}}{2 L D^{2}} .
\end{aligned}
\tag{19}
\]
\textbf{[Step 1]} The cancellation follows from algebra:
\[
-\frac{\mathcal G^{2}}{L D^{2}}+\frac{1}{2}\frac{\mathcal G^{2}}{L D^{2}}
= -\frac{1}{2}\frac{\mathcal G^{2}}{L D^{2}} .
\]

\paragraph{Case B: $\displaystyle \gamma_{k}=1$}
In this case the definition of $\gamma_{k}$ implies
$\mathcal G(x_{k})\ge L D^{2}$.  Using $\gamma_{k}=1$ in \eqref{17} we obtain
\[
f(x_{k+1}) \le f(x_{k}) - \mathcal G(x_{k}) + \frac{L D^{2}}{2}.
\tag{20}
\]
Since $\mathcal G(x_{k})\ge L D^{2}$,
\[
-\mathcal G(x_{k}) + \frac{L D^{2}}{2}
\le -\frac{\mathcal G(x_{k})}{2}.
\tag{21}
\]
Thus,
\[
f(x_{k+1}) \le f(x_{k}) - \frac{\mathcal G(x_{k})}{2}
           \le f(x_{k}) - \frac{\mathcal G(x_{k})^{2}}{2 L D^{2}},
\tag{22}
\]
where the last inequality uses $\mathcal G(x_{k})\ge L D^{2}$ again.

\subsection*{Step 7 – Unified descent inequality}
Both cases lead to the same bound:
\[
\boxed{\,f(x_{k+1}) \le f(x_{k}) - \frac{\mathcal G(x_{k})^{2}}{2 L D^{2}}\,}.
\tag{23}
\]
\textbf{[Step 2]} This inequality is the cornerstone of the convergence
analysis.

\subsection*{Step 8 – Telescoping sum}
Sum \eqref{23} over $k=0,\dots,T-1$:
\[
\sum_{k=0}^{T-1}\frac{\mathcal G(x_{k})^{2}}{2 L D^{2}}
   \le f(x_{0}) - f(x_{T})
   \le f(x_{0}) - f^{*}.
\tag{24}
\]
Rearranging,
\[
\frac{1}{T}\sum_{k=0}^{T-1}\mathcal G(x_{k})^{2}
   \le \frac{2 L D^{2}\bigl(f(x_{0})-f^{*}\bigr)}{T}.
\tag{25}
\]

\subsection*{Step 9 – From average of squares to minimum gap}
Since the minimum of a set of non‑negative numbers is not larger than the
average of their squares,
\[
\min_{0\le k\le T-1}\mathcal G(x_{k})
   \;\le\;
\sqrt{\frac{1}{T}\sum_{k=0}^{T-1}\mathcal G(x_{k})^{2}}.
\tag{26}
\]
Combining \eqref{25} and \eqref{26} yields exactly the bound \eqref{8}:
\[
\min_{0\le k\le T-1}\mathcal G(x_{k})
   \le
\sqrt{\frac{2 L D^{2}\bigl(f(x_{0})-f^{*}\bigr)}{T}} .
\tag{27}
\]

Thus Theorem~\ref{thm:main} is proved.  $\square$

%%%%%%%%%%%%%%%%%%%%%%%%%%%%%%%%%%%%%%%%%%%%%%%%%%%%%%%%%%%%
\section{Rate Derivation (With Constants)}
\label{sec:rate}
The constant appearing in \eqref{8} is
\[
C \;:=\; \sqrt{2 L D^{2}\bigl(f(x_{0})-f^{*}\bigr)} .
\]
Therefore the algorithm enjoys an $O(1/\sqrt{T})$ decay of the Frank–Wolfe
gap.  If one uses the diminishing stepsize $\gamma_{k}=2/(k+2)$ instead of
\eqref{7}, a slightly weaker bound
\[
\min_{0\le k\le T-1}\mathcal G(x_{k}) \le
\frac{2 L D^{2}}{T}
\bigl(f(x_{0})-f^{*}\bigr)
\tag{28}
\]
can be derived by the same technique (the algebra is omitted for brevity).

%%%%%%%%%%%%%%%%%%%%%%%%%%%%%%%%%%%%%%%%%%%%%%%%%%%%%%%%%%%%
\section{Special Cases}
\label{sec:special}
\subsection*{Convex $f$}
When $f$ is convex, the FW gap upper‑bounds the primal optimality gap:
\[
\mathcal G(x) \;\ge\; f(x)-f^{*}.
\]
Consequently, \eqref{8} implies the classic $O(1/\sqrt{T})$ bound on the
optimality error for the convex case with the adaptive stepsize, and the
$O(1/T)$ bound when $\gamma_{k}=2/(k+2)$.

\subsection*{Polyhedral feasible set}
If $\mathcal C$ is a polytope, the LMO reduces to solving a linear program,
which can be done in $O(d)$ time for many structured polytopes (e.g. the
probability simplex).  The convergence rate remains unchanged.

\subsection*{Strongly smooth but non‑convex}
If, in addition to $L$‑smoothness, the Hessian satisfies
$\|\nabla^{2}f(x)\|\le L$ (i.e. the function is \emph{quadratically}
smooth), the same proof goes through unchanged; the constant $L$ can be
replaced by the tighter quadratic smoothness constant, potentially improving
the practical rate.

%%%%%%%%%%%%%%%%%%%%%%%%%%%%%%%%%%%%%%%%%%%%%%%%%%%%%%%%%%%%
\section*{References}
\begin{enumerate}
\item J. Jaggi, “Revisiting Frank–Wolfe: Projection‑Free Sparse Convex Optimization,” \emph{ICML}, 2013.
\item H. L. Bauschke, P. L. Combettes, \emph{Convex Analysis and Monotone Operator Theory in Hilbert Spaces}, 2011.
\item A. L. (2022). “Non‑convex Frank–Wolfe with Adaptive Stepsizes,” \emph{Optimization Letters}.
\end{enumerate}

\end{document}